This chapter presents the recommender core in its current state of development. Consistent with
the progress reported to the supervisor, the retrieval stage is currently being implemented; the
remaining pipeline stages are planned as next steps and have not yet been built.

\section{Overview}
\label{sec:rec-overview}

The recommender core is designed to combine three inputs --- the current question, the user's own
past queries, and their behavioral profile (Chapter~\ref{ch:behavior}) --- into three follow-up
query suggestions. At the current stage, retrieval of semantically similar past queries is under
implementation; candidate generation, schema-grounding validation, and ranking are planned as the
next steps of this work and are not yet built.

\section{Retrieval (In Progress)}
\label{sec:rec-retrieval}

Embedding-based retrieval is currently being implemented. Sentence embeddings and cosine
similarity will be used to find previous user queries that are semantically related to the
current question, even when they use different words. This retrieval step is intended to supply
the similar-past-queries context used by the candidate-generation stage described below.

\section{Planned Next Steps}
\label{sec:rec-generation}
\label{sec:rec-validation}
\label{sec:rec-ranking}

The next step is to combine the current question, the retrieved similar past queries, and the
behavioral profile (Chapter~\ref{ch:behavior}). An LLM will generate candidate follow-up
questions from these three inputs; the candidates will then be validated against the domain
schema, ranked, and diversified to select the best three. These stages are not yet implemented.
Ollama is being used as a local LLM environment for development and testing.
